% =========================================================================
% CAPITOLO 10 — DISCUSSIONE
% =========================================================================

\chapter{Discussione}
\label{ch:discussion}

\noindent
{\rmfamily\itshape
\begin{tabular}[t]{@{}l@{\hspace{0.3cm}}p{12.5cm}@{}}
\textls[50]{\textbf{Trattati:}} &
\textls[50]{Interpretazione · Aspettative riviste · Meccanismo e condizioni al contorno ·
Implicazioni teoriche, pratiche e di sostenibilità · Ciò che lo studio può e non può affermare}
\end{tabular}
}

\vspace{0.5cm}

Il Capitolo~\ref{ch:results} ha riportato ciò che mostrano le risposte codificate; questo capitolo chiede cosa significhino. Legge i risultati rispetto alle aspettative che hanno inquadrato lo studio con utenti, e indica dove l'evidenza sostiene la tesi, dove la complica e dove non può pronunciarsi. Con venticinque partecipanti, lo studio non può dimostrare un'affermazione generale; ciò che può fare è mostrare se i pattern che è stato costruito per cercare compaiano davvero --- e il risultato più utile di un primo studio spesso non è una conferma ma una comprensione più affinata di dove l'approccio aiuta e perché.

% -------------------------------------------------------------------------
\section{Rileggere le aspettative alla luce dello studio realizzato}
\label{sec:disc-translate}

Le aspettative per la valutazione, E3.1 ed E3.2, sono state scritte presto, per un disegno che poi è cambiato: tre condizioni di crescente ricchezza esperienziale e un protocollo think-aloud sono diventate due condizioni e un questionario scritto auto-somministrato (Capitolo~\ref{ch:userstudy}). Vanno quindi lette in forma tradotta. E3.1 --- che una maggiore ricchezza esperienziale si sarebbe mostrata in \emph{come} i partecipanti articolano la propria comprensione anziché in un punteggio di comprensione --- è letta come un confronto tra le due condizioni. E3.2 --- che le codifiche percettive avrebbero prodotto interpretazioni che il protocollo poteva far emergere, corrispondenti o meno alla lettura intesa --- è letta come riferita a ciò che le risposte scritte rivelano.

Questa traduzione è un limite, ed è enunciata come tale; ma non indebolisce le aspettative, perché E3.1 era prudente esattamente nel modo che si è rivelato importante. Non prevedeva che la condizione esperienziale aumentasse la comprensione; prevedeva una differenza nell'articolazione, non nel punteggio --- che è precisamente ciò che i risultati mostrano.

% -------------------------------------------------------------------------
\section{E3.1: l'aspettativa prudente regge, quella ingenua no}
\label{sec:disc-e31}

Il singolo risultato più netto è che la condizione esperienziale non ha prodotto un guadagno uniforme di comprensione. Nelle tre scene la direzione della differenza non era la stessa (Figura~\ref{fig:grasped}): pari nella Scena~I, a favore della condizione sperimentale nella Scena~II, a favore del controllo nella Scena~III. Uno studio che si fosse aspettato che la presentazione esperienziale alzasse ovunque la comprensione sarebbe messo in imbarazzo da questo pattern. Il presente studio no, perché non se lo è mai aspettato.

Se l'aspettativa fosse stata quella ingenua --- che una rappresentazione più ricca e in movimento produca semplicemente una comprensione migliore --- i dati l'avrebbero rifiutata. Ciò che è cambiato con la ricchezza esperienziale non è stato il punteggio di comprensione ma il carattere della comprensione, e dove essa fosse affatto raggiungibile: esattamente la forma di E3.1. L'avvertimento che le visualizzazioni animate possono apparire efficaci solo perché introducono di nascosto informazione che la versione statica non ha \parencite{tversky2002} è il motivo per cui le due condizioni sono state allineate per provenienza, inquadramento e unità di misura, con la versione esperienziale che aggiunge solo la dimensione temporale che lo studio si propone di esaminare. Il confronto piatto nella Scena~I è onesto anziché deludente.

% -------------------------------------------------------------------------
\section{Dove vive la differenza, e perché}
\label{sec:disc-mechanism}

Se l'effetto non è un aumento uniforme, la domanda interessante è dove compaia e perché.

La Scena~II riguarda la ricorrenza: se il sistema ripassi per stati che ha già occupato. Uno scatter immobile non può mostrarlo, perché senza un ordinamento nel tempo non c'è nulla che riveli un ritorno. I risultati corrispondono esattamente. Nessun partecipante di controllo è stato codificato come avente colto la ricorrenza, contro circa un quarto del gruppo sperimentale, con in aumento anche la comprensione parziale (Capitolo~\ref{ch:results}). Il meccanismo non è dedotto; lo enunciano i partecipanti stessi. Ogni risposta \textsc{Interaction-aided} --- che attribuisce la comprensione al mettere in pausa, riavvolgere o isolare parte del display --- veniva dalla condizione sperimentale, e ogni risposta \textsc{Static-limit-named} --- che nomina l'assenza di un ordinamento temporale come la ragione per cui la ricorrenza non poteva essere giudicata --- veniva dal controllo (Figura~\ref{fig:emergent}): le due facce della stessa affermazione, dette da lati opposti del disegno.

Un terzo codice emergente va contro la condizione esperienziale, ed è riportato qui per la stessa ragione per cui lo è l'inversione della Scena~III. \textsc{Anomaly-confusion}, che segna la difficoltà con la codifica usuale/insolito, è comparso cinque volte nella condizione sperimentale contro una nel controllo (\S\ref{sec:res-emergent}). Il principio~P3 assegna l'anomalia a texture e agitazione, sul presupposto che l'irregolarità attiri l'attenzione senza una legenda \parencite{bertin1983, ware2021}. Queste risposte suggeriscono che attirare l'attenzione non è lo stesso che trasmettere significato: i partecipanti notavano il disturbo e non sapevano dire cosa significasse. La salienza pre-attentiva pare sufficiente per il rilevamento e insufficiente per l'interpretazione --- un limite di P3 così come enunciato, e un obiettivo specifico di revisione.

La Scena~III è la prova che rende tutto questo credibile, perché va nell'altro verso. Qui la condizione di controllo ha fatto meglio, raggiungendo ogni partecipante contro circa tre quarti del gruppo sperimentale. La ragione è istruttiva anziché imbarazzante. La Scena~III chiede se l'accordo tra tre modelli cambi nel tempo, e una figura statica mostra l'intera serie temporale in una volta, dispiegando la convergenza da leggere in un'unica vista. La versione esperienziale chiede al partecipante di costruire quella stessa lettura muovendosi tra i controlli, e alcuni non l'hanno completata. Una seconda spiegazione non può essere separata dalla prima con questo disegno: la Scena~III veniva per ultima per ogni partecipante, e la versione esperienziale chiede attenzione sostenuta proprio dove l'attenzione è meno probabile che resti. L'ordine fisso delle scene (\S\ref{sec:eval-limitations}) fa sì che i due resoconti siano confusi, e lo studio non può decidere tra essi. Sotto entrambe le letture, la lezione non è che la forma esperienziale abbia fallito. È che la presentazione esperienziale non è gratuita. Aiuta di più dove la struttura è invisibile in una figura immobile, come nella Scena~II, e può costare attenzione dove una figura statica già la dispiega con chiarezza, come nella Scena~III. Un approccio che ammette questo è più forte di uno che rivendica un vantaggio uniforme, perché può dire \emph{dove} si applica.

% -------------------------------------------------------------------------
\section{E3.2 e l'errore sulla funzione di corrente}
\label{sec:disc-e32}

Il risultato più sostanziale dello studio è una lettura errata, ed E3.2 anticipava esattamente questo genere di cosa: che i non specialisti avrebbero interpretato la rappresentazione in modi che valeva la pena far emergere, che l'interpretazione corrispondesse o no a quella intesa. Quattordici partecipanti su venticinque hanno letto l'altezza della curva della funzione di corrente nella Scena~I come la velocità dell'acqua, anziché come il trasporto cumulato che rappresenta (Capitolo~\ref{ch:results}). La curva invita quella lettura. La sua altezza è la cosa più saliente sul display, e l'altezza si legge naturalmente come quantità o velocità; la lettura corretta, che direzione e intensità vivono nella pendenza della curva, richiede un'inferenza che un non specialista non fa spontaneamente.

Questa lettura errata non è specifica di un dominio, e il Capitolo~\ref{ch:literature} la nomina già. La funzione di corrente è un integrale cumulato del trasporto, la sua pendenza il tasso locale (Equazione~\ref{eq:velocity}); leggere la sua altezza come velocità è la confusione stock--flusso che \textcite{sterman2008} e \textcite{cronin2009} documentano in adulti matematicamente competenti. Lo studio aggiunge che la confusione sopravvive alla traduzione in un dominio e in un idioma diversi: i partecipanti hanno incontrato una curva oceanografica anziché un diagramma stock-e-flusso e hanno fatto la stessa sostituzione --- un limite noto dell'architettura cognitiva che riappare esattamente dove il filone delle scienze cognitive prevedeva, trasformando il risultato della Scena~I da anomalia in conferma.

Due ulteriori pattern trasformano questo da errore in risultato. Il primo è che l'errore era più frequente nella condizione sperimentale, non meno (Figura~\ref{fig:misconception}). A prima vista questo contraddice la tesi. Non è così, una volta affiancato il secondo pattern: l'errore calava nettamente con la familiarità con le visualizzazioni di dati, da circa tre quarti dei partecipanti meno familiari a un terzo dei più familiari. Il display esperienziale mostra particelle in movimento, e le particelle in movimento rendono la lettura di velocità più immediata, non meno. La ricchezza percettiva che ha aiutato i partecipanti a vedere la ricorrenza nella Scena~II qui lavora contro di loro, rendendo la lettura seducente più facile proprio per chi ha meno strumenti per correggerla. Questo non è tanto un fallimento dell'approccio quanto una lezione su di esso: la salienza percettiva è potente in entrambe le direzioni, e una rappresentazione che rende vivida una struttura può rendere vivida anche una lettura errata adiacente.

Che le due condizioni differissero anche nella familiarità media con le visualizzazioni di dati significa che condizione e familiarità non possono essere pienamente separate a questa dimensione campionaria (Capitolo~\ref{ch:results}), perciò la lettura congiunta è esplorativa. La direzione è comunque abbastanza chiara da portare una conseguenza di design, ripresa nel Capitolo~\ref{ch:conclusion}: un'iterazione futura deve rendere percepibile la natura cumulata della curva anziché lasciare che sia l'osservatore a fornirla.

L'errore dà anche a E3.2 il suo spigolo di sicurezza-ma-errore. \textcite{fischer2020} descrivono casi in cui un'alta fiducia si attacca a una lettura errata, e la Scena~I ne contiene: quattro partecipanti hanno valutato la propria fiducia a sette o più pur leggendo la curva solo parzialmente, ciascuno portando l'errore sulla funzione di corrente (Figura~\ref{fig:confidence}). La fiducia, in questa scena, non segue la comprensione. Le persone più sicure della propria lettura erano tra quelle che avevano letto male lo strumento, che è precisamente il motivo per cui lo studio ha tenuto separate comprensione percepita e dimostrata anziché lasciare che una valutazione di fiducia facesse le veci di una presa.

% -------------------------------------------------------------------------
\section{Ciò che lo studio non può affermare}
\label{sec:disc-limits}

Tre limiti vincolano queste letture, e nominarli è parte dell'argomentazione anziché una scusa aggiunta a essa.

Il campione è piccolo ed è stato reclutato per convenienza, perciò nulla qui si generalizza a una popolazione e i test esplorativi non portano peso inferenziale. Il confronto della Scena~III, dove il controllo ha fatto meglio, è reale quanto quello della Scena~II, e lo studio riporta entrambi anziché il solo lusinghiero. Il confondimento tra condizione e familiarità con le visualizzazioni di dati significa che i due pannelli del risultato sull'errore non sono indipendenti: il risultato è una direzione da testare, non un effetto accertato. E la codifica poggia su un singolo ricercatore supportato da un primo passaggio automatico, perciò non può essere riportata alcuna affidabilità tra codificatori (Capitolo~\ref{ch:userstudy}). Nessuno è fatale a un primo studio; ciascuno è una ragione perché il prossimo sia più ampio, bilanciato sulla familiarità pregressa, codificato in modo indipendente e controbilanciato sull'ordine delle scene.

% -------------------------------------------------------------------------
\section{Implicazioni}
\label{sec:disc-implications}

\paragraph{Teoriche.} Lo studio fornisce un caso in cui la confusione stock--flusso di \textcite{sterman2008} riappare in un idioma visivo scientifico anziché in un diagramma da manuale (\S\ref{sec:disc-e32}), e uno in cui il principio di equivalenza di \textcite{tversky2002} fa un lavoro reale: con provenienza, inquadramento e unità mantenuti costanti e solo la dimensione temporale aggiunta, il vantaggio dell'animazione non compare in modo uniforme. Qualifica anche la letteratura sull'incertezza animata \parencite{hullman2015hops, kale2019hops}: la Scena~III ha animato l'incertezza e non ha superato la baseline statica, il che suggerisce che il beneficio possa dipendere da se l'alternativa statica offra già una vista dell'intera serie.

\paragraph{Pratiche.} Per chiunque progetti una rappresentazione di un sistema invisibile, il risultato è una regola di selezione, non un avallo: il movimento merita il suo costo dove la struttura è temporale e non disponibile in una figura immobile, e aggiunge carico attentivo senza comprensione dove una figura statica la dispiega già. La salienza va verificata in entrambe le direzioni: il campo di particelle che ha reso visibile il contro-flusso nella Scena~I ha reso anche più disponibile la lettura errata di velocità, soprattutto per chi era meno attrezzato a correggerla.

\paragraph{Sostenibilità e comunicazione.} Che la comprensione non segua la chiarezza percepita estende \textcite{hegarty2012} a un formato esperienziale, e riguarda come il materiale sul clima è preparato per il pubblico non specialista: un artefatto valutato sulla sola preferenza o sul solo coinvolgimento può essere ottimizzato nella direzione sbagliata. La memoria di lavoro è il vincolo stringente \parencite{sweller1988}, e una rappresentazione che la spende in decorazione anziché in struttura la spende male.

% -------------------------------------------------------------------------
\section{Cosa significa per H-A e H-B}
\label{sec:disc-umbrella}

Le ipotesi ombrello hanno inquadrato l'intera indagine: che le rappresentazioni convenzionali siano percettivamente subottimali per i non specialisti perché trattano la complessità come un oggetto da leggere anziché una dinamica da abitare (H-A), e che le rappresentazioni esperienziali derivate dal machine learning possano sostenere la comprensione rendendo la struttura percettivamente disponibile (H-B). Venticinque partecipanti non possono dimostrare nessuna delle due. Entrambe trovano il loro sostegno più netto nella Scena~II, dove la rappresentazione statica ha lasciato la ricorrenza semplicemente irraggiungibile e i partecipanti l'hanno detto con parole proprie, e in coloro che hanno attribuito la propria comprensione all'agire sul display.

Il contributo più prezioso è il confine che lo studio traccia attorno a H-B, esposto nella \S\ref{sec:disc-mechanism} e nella \S\ref{sec:disc-implications}: la presentazione esperienziale non è uniformemente migliore, e il suo potere percettivo può approfondire una lettura errata con la stessa facilità di una corretta. La tesi emerge non come dimostrata ma come affinata --- la domanda non è più se la rappresentazione esperienziale aiuti, ma per quali strutture, per quali osservatori e a quale costo percettivo.
